\chapter*{实验总则}
\addcontentsline{toc}{chapter}{实验总则}

实验的目标是验证一个明确机制。开始前应记录问题、输入、对照条件和可证伪判据。若只能运行原理计算部分，应如实界定范围；不要把未运行的模型迁移、GPU内核或多机部署写为通过。

\section*{环境及资源}
为实验使用独立的Python环境，按Notebook的依赖单元准备库并记录实际版本。先检查数据与模型制品的来源、修订、许可及下载规模，再决定是否运行相应单元。CPU原理计算、GPU训练和真实服务测试分别计量，不能按Notebook文件名假设资源需求。记录设备、可用内存、精度、批次、序列长度和预算；不具备所需设备时，保留该项为未验证。

\section*{执行顺序及状态}
先阅读整个实验的数据流和资产流，再从初始化单元按依赖顺序执行。修改数据、种子或模型配置后应重新检查下游状态；最终提交前从干净内核重放所声明的执行范围。保存输入制品身份和输出记录，避免隐藏变量使结果只能在当前会话中复现。需要联网下载、较长训练或外部服务的部分应单独列出条件和成本。

\section*{数值及性能验证}
数值对照使用相同输入、参数、掩码和目标归一化方式；报告最大误差或相对误差，并说明容差依据。随机实验应记录种子和重复次数。性能测量需区分初始化、预热与稳定运行；加速器计时考虑同步，报告输入长度分布与统计量。通过形状检查不意味着数值正确，通过数值对照也不意味着获得性能提升。

\section*{提交材料}
每份报告至少包含实验目标、环境与资产、操作范围、预期判据、实际证据和失败解释。标准记录格式与各任务专用字段见附录\ref{app:experiment-reporting}《实验报告模板》。图表应可追溯至原始结果；保留错误样本及未通过检查，不只展示成功案例。若使用构造数据或状态演算，应明确其只能验证算法关系，不能替代真实模型质量、服务吞吐或业务收益。

\section*{从理论到实现}
先在教材中确定定义与前提，再在本册中完成推导练习，随后在Notebook中定位对应对象。原理实现负责说明机制，生产库负责完整接口与执行路径；迁移时比较的是契约、数值和边界，不能仅因两段代码都能运行便认定它们等价。


\section*{框架实践路线}
微调框架实践见实验\ref{lab:42_peft}，先完成数据、监督位置、训练和适配器重载，再对照集成训练框架。分布式实践见实验\ref{lab:A40_training_optimization}，保持训练目标与有效批量，研究状态分片、检查点和恢复。推理部署实践见实验\ref{lab:60_inference_deployment}，先完成vLLM服务路径，再迁移到SGLang并比较质量、负载和故障行为。

三条路线以本册中的命令、配置和配套训练入口为操作材料。版本号用于定位接口核查基线，完整环境、硬件支持与性能结论须由实际运行建立。选择框架时先确定职责和验收条件，同类工具的迁移应复用相同数据、模型身份和评价口径。
